\begin{trapbox}

	\begin{description}[style=nextline, leftmargin=2.8em, labelsep=0.5em, itemsep=0.35em]
		\item[\textbf{\textcolor{CTrap}{易错点一（忽略判别式条件）}}]
			未验证\(\Delta'\)是否大于0，就直接使用极值乘积符号判定零点个数。
		\item[\textbf{\textcolor{CTrap}{正确理解（先验条件）：}}]
			使用乘积判别前必须首先计算\(\Delta'\)，只有当\(\Delta'>0\)（即函数有两个极值点）时，该判定才适用。若\(\Delta'\leq0\)，函数单调，直接得零点唯一。
		\item[\textbf{\textcolor{CTrap}{错因分析（机械套用）：}}]
			死记硬背乘积符号结论，忽略其前提条件，导致在单调函数上错误地认为极值存在。
	\end{description}

	\medskip
	\begin{description}[style=nextline, leftmargin=2.8em, labelsep=0.5em, itemsep=0.35em]
		\item[\textbf{\textcolor{CTrap}{易错点二（乘积符号记反）}}]
			误认为乘积为正时三个零点，乘积为负时一个零点。
		\item[\textbf{\textcolor{CTrap}{正确理解（记忆口诀）：}}]
			乘积负三零，正一零，零二重。即\(f(x_1)f(x_2)<0\)时三个零点，\(>0\)时一个零点，\(=0\)时两个零点（重根）。
		\item[\textbf{\textcolor{CTrap}{错因分析（图像不清）：}}]
			对三次函数图像缺乏直观理解，混淆了极值符号与零点个数的对应关系。
	\end{description}

	\medskip
	\begin{description}[style=nextline, leftmargin=2.8em, labelsep=0.5em, itemsep=0.35em]
		\item[\textbf{\textcolor{CTrap}{易错点三（重根计数遗漏）}}]
			当乘积为零时，只算出一个零点，忽略了重根计为零点的事实。
		\item[\textbf{\textcolor{CTrap}{正确理解（重根计数）：}}]
			乘积为零意味着函数有一个极值点恰好是零点（即重根），同时还有另一个零点，因此总零点数为2。（注意：重根次数计入零点个数。）
		\item[\textbf{\textcolor{CTrap}{错因分析（概念不清）：}}]
			对“重根”理解为“只有一个根”，忘记重根在计数时按重数计算；或误认为极值点处为零不算零点。
	\end{description}

\end{trapbox}
